\subsection{Terrain pipeline and sites}\label{sec:terrain}
LunarSim builds a landing tile from a real coarse elevation source plus synthetic layers. The coarse layer is a LROC NAC digital terrain model; we use the Apollo~11 and Apollo~16 products (about 2~m per pixel). Crater, rock and fine-detail layers are added on top because no raster resolves them. Fine detail is injected as a high-pass component weighted by a Gaussian region-of-interest mask, so the coarse layer keeps the low frequencies. A separate high-resolution patch can be generated around a point of interest. Tiles are pre-generated into pools (a training half and a held-out half per site); all experiments here draw from the held-out pools. \todo{state pool sizes and the longitude ranges after P0 confirms them}

\subsection{Landing criterion}\label{sec:criterion}
A landing counts as safe when touchdown speed, tilt, body rate and footpad flatness are inside their limits: vertical speed $\leq 1.0$~m/s, horizontal speed $\leq 1.2$~m/s, tilt $\leq 15^{\circ}$, body rate $\leq 0.5$~rad/s, and an effective footpad term $\leq 0.16$~m. The default footpad term (\texttt{plane}) folds two checks into one number: the terrain slope under the four pads ($\leq 10^{\circ}$ by default) and the roughness, i.e.\ the pad-height residual after removing the best-fit plane. A legacy definition (\texttt{spread}) uses the raw maximum-minus-minimum of the four pad heights. Every result file stores the active criterion so that numbers are never compared across definitions.

\subsection{Isaac Sim and Isaac Lab backend}\label{sec:isaac}
The landing environment runs in Isaac Sim 6.0.1 with Isaac Lab on an NVIDIA GPU. The terrain is added to the USD stage as a collision surface for each environment, and the ground height used for contact and for the footpad term is read from the actual collision mesh, not from the source array. Physics runs at 120~Hz and the controller acts every 0.05~s. The vehicle is a rigid-body model with Apollo Lunar Module mass properties (dry mass 6938~kg, descent propellant 8165~kg) and attitude control through reaction-control torques.

\subsection{Stage and behavioural probe}\label{sec:stage}
All experiments use the stage \texttt{orbit\_descent}: the vehicle is released at 200~m altitude within 80~m of the target on a 1680~m tile, with 10--30~m/s horizontal speed, and the episode lasts at most 60~s. The probe is a ZEM/ZEV guidance controller (zero-effort-miss / zero-effort-velocity) with fixed gains. It is deterministic, so for a given terrain and release state the only sources of difference between runs are the simulator setting under test and the physics engine's own nondeterminism, which Experiment~E0 measures directly.
